\section{The complete cube curl and a failed literal subtraction}
\label{sec:r11-cube-curl-extraction}

\paragraph{Status and scope.}
The finite curl formulas below are \textbf{PROVED THIS ROUND}.
The proposed literal cancellation of the GR remainder's first curl
is a \textbf{ROUTE FAILURE}: an explicit native rank-four facet
violates it.  The construction contributors to this module are
\texttt{/root/r11\_kernel} and \texttt{/root}; neither is eligible
as a nonconstructor reviewer.  The formula was derived before the
facet calculation, and the latter was also checked independently
by two finite exterior-algebra calculations.  Those calculations
are calibrations, not replacements for the carrier proof below.

The conclusion does not exclude an actual $D$-primitive for the
nonzero curl, a different extending triple, or a relative cube
witness.  In particular it is neither a Strong G counterexample
nor an entire actual OPEN-KERNEL counterexample.

\subsection{Notation and the exact extraction}

Keep the eight original actual norm potentials and the six oriented
faces of GR12, in rank $r\ge3$.  Put
\[
 s=r-1,\qquad d=r-2,\qquad \ell=r-3=d-1,\qquad
 \theta_v=\Phi_v-\Phi_0,\qquad C_0=R-D\Phi_0,
 \qquad C_v=C_0-D\theta_v.
 \tag{CF1}
\]
For face $\nu$, put $\kappa_\nu=\Psi_\nu-\Phi_0$ with its
specified original CS cap.  Its ordered weights are
$(a_\nu,b_\nu)$; an oriented edge $e=(i,j)$ has signed weight
$c_e$.  The face signs $\sigma_\nu$ are exactly GR12.  All these
objects live on the whole original punctured open $W$.

For a numerical parameter $w\in[0,1]$ set
\[
 R_{\nu,w}=C_0-wD\kappa_\nu,\qquad
 p_{v,\nu}(w)=\theta_v-w\kappa_\nu.
 \tag{CF2}
\]
These are finite polynomial expressions in actual forms, not
assumed to be actual tuple norm presentations for intermediate
$w$.  Write $h_j$ for the complete homogeneous polynomial and
$\partial_Rh_j(R,\ldots)$ for its formal derivative in the first
displayed variable.  The other curvature arguments are held fixed.

Define the following completely specified actual form:
\[
\begin{split}
 \mathcal J_{\nu,d}(w)
  ={}&\sum_{e=(i,j)\in\nu}\Bigl[
       d''\theta_i\wedge d''\theta_j\,
          \partial_Rh_d(R_{\nu,w},C_i,C_j)\\
 &\hspace{15mm}+d''(c_e)_+\wedge d''\theta_j\,
          \partial_Rh_d(R_{\nu,w},C_j)\\
 &\hspace{15mm}+d''(c_e)_-\wedge d''\theta_i\,
          \partial_Rh_d(R_{\nu,w},C_i)\Bigr]\\
 &+d\,d''a_\nu\wedge d''b_\nu\,R_{\nu,w}^{d-1}.
\end{split}
 \tag{CF3}
\]
Here $(c_e)_+=\max(c_e,0)$ and $(c_e)_-=\min(c_e,0)$ are
the full actual piecewise affine unit-weight functions; no sign
wall is discarded.

\begin{proposition}[Full first curl of the metric-independent GR triple]
For the exact $E_{\rm glob}$ of GR15, its first component satisfies
the literal actual identity
\[
 \boxed{\quad
 d''(E_{\rm glob})_A
   =\sum_{\nu=1}^6\sigma_\nu\,
      d''\kappa_\nu\wedge\int_0^1\mathcal J_{\nu,d}(w)\,dw.
 \quad}
 \tag{CF4}
\]
The integration consists only of integrating finitely many real
polynomial coefficients.  It introduces no completion, current
primitive or analytic parameter pushforward.
\end{proposition}

\begin{proof}
Use GR's finite two-colour algebra and put $w(t)=1-\rho(t)$,
where $w(0)=1$, $w(1)=0$.  The full reference curvature is
\[
 \widetilde R_\nu
  =R_{\nu,w}-\epsilon'\dot w\,d''\kappa_\nu
       +\epsilon''\dot w\,d'\kappa_\nu
       -\epsilon'\epsilon''\ddot w\,\kappa_\nu.
 \tag{CF5}
\]
The projective curvatures $C_v$ are independent of $t$ and of
both exterior parameter variables, since
$\widetilde R_\nu-\widetilde Dp_{v,\nu}=C_v$.

For actual scalars write
$s''(u,v)=u\,d''v-v\,d''u$.  Its formal counterpart has no
single $\epsilon'$ coefficient.  Moving $\epsilon'$ to the
left through the odd symbol $s''$ introduces a minus sign,
which cancels the minus sign of the first parameter coefficient
in (CF5).  Consequently the \emph{complete} $F_1^+$ coefficient
in GR11 for this face is
\[
\begin{split}
 (F_\nu)_1^+
 =\frac{\dot w}{2}\Bigl\{\sum_{e=(i,j)}\bigl[
 &s''(p_i(w),p_j(w))\wedge d''\kappa_\nu\,
                         \partial_Rh_d(R_{\nu,w},C_i,C_j)\\
 +{}&s''((c_e)_+,p_j(w))\wedge d''\kappa_\nu\,
                         \partial_Rh_d(R_{\nu,w},C_j)\\
 +{}&s''((c_e)_-,p_i(w))\wedge d''\kappa_\nu\,
                         \partial_Rh_d(R_{\nu,w},C_i)\bigr]\\
 +{}&d\,s''(a_\nu,b_\nu)\wedge d''\kappa_\nu\,
                                      R_{\nu,w}^{d-1}\Bigr\}.
\end{split}
 \tag{CF6}
\]
This extraction does not declare the other single or double
parameter coefficients zero.  They remain present in GR5--GR9;
the first component of GR15 simply uses this particular one.

The first component in GR15 is
$-\int_0^1\sum_\nu\sigma_\nu(F_\nu)_1^+\,dt$.
Changing variables from $t$ to $w$ reverses the endpoints and
cancels this minus sign.  For each fixed numerical $w$, all
curvature factors in (CF6) are $d''$-closed, and
\[
 d''\bigl(s''(u,v)\wedge d''\kappa\bigr)
                  =2d''u\wedge d''v\wedge d''\kappa.
 \tag{CF7}
\]
The factor two cancels the one half in (CF6).  In every resulting
wedge, replacing $d''p_v(w)=d''\theta_v-wd''\kappa$ by
$d''\theta_v$ makes no change, because another $d''\kappa$
is already present.  Moving that last one-form to the front
passes two odd factors and has positive sign.  The result is
exactly (CF3)--(CF4).
\end{proof}

In rank three $d=1$, all three formal derivatives in (CF3)
are one, as is the final factor $dR_{\nu,w}^{d-1}$.  In rank
four $d=2$, the four integrated coefficients, in their displayed
order, are
\[
 \begin{gathered}
 C_0+C_{\Psi_\nu}+C_i+C_j,\qquad
 C_0+C_{\Psi_\nu}+C_j,\qquad
 C_0+C_{\Psi_\nu}+C_i,\qquad C_0+C_{\Psi_\nu},\\
             C_{\Psi_\nu}=C_0-D\kappa_\nu.
 \end{gathered}
 \tag{CF8}
\]
This follows by integrating $2R_{\nu,w}$ and using
$2C_0-D\kappa_\nu=C_0+C_{\Psi_\nu}$.  No curvature factor
has been cancelled in an actual quotient.

\subsection{A separate formula for the literal raw-face curl}

The raw cube and its metric-independent remainder need not have
the same first curl.  For comparison, let a face have unshifted
potentials $\Phi_i$, common reference $R$, and actual normalizer
$\Psi$.  Set $p_i=\Phi_i-\Psi$, $R_*=R-D\Psi$ and
$C_i=R-D\Phi_i$.  For the \emph{entire} CS36 expression
$q[\Psi]$, including full MT and CPT, one has
\[
\begin{split}
 d''q[\Psi]_A
 ={}&\frac12\sum_{e=(i,j)}d''|c_e|\wedge d''p_i\wedge d''p_j\,
                                      h_\ell(R_*,C_i,C_j)\\
 &+d''\Psi\wedge\sum_{e=(i,j)}\Bigl[
      d''\Phi_i\wedge d''\Phi_j\,
                                      h_\ell(R_*,R,C_i,C_j)\\
 &\hspace{19mm}+\tfrac12 d''c_e\wedge d''\Phi_i\,
                                      h_\ell(R_*,R,C_i)\\
 &\hspace{19mm}+\tfrac12 d''c_e\wedge d''\Phi_j\,
                                      h_\ell(R_*,R,C_j)\Bigr]\\
 &+d''\Psi\wedge d''a_\nu\wedge d''b_\nu\,
                                      h_\ell(R_*,R).
\end{split}
 \tag{CF9}
\]

\begin{proof}
The first component of the scalar triple in CS6 has no
contribution.  Differentiating the numerator of CS9 in the
second colour gives three times
$d''|c_e|\wedge d''p_i\wedge d''p_j$.
The one-third factor of CS9 and the one-half factor of the full
edge triple therefore give the first line of (CF9), initially
with $R_*-D|c_e|$ in place of its first curvature argument.
Every difference term contains $D|c_e|$.  Its carrier is on
$c_e=0$, where the tangential first derivative of $|c_e|$ is
zero.  The explicit factor $d''|c_e|$ therefore kills every
such term, on all refinements.  This proves the first line
with its stated curvature polynomial, without cancelling a
curvature factor.

For MT's complete coefficient pair $f_k$, direct differentiation
of MT6 gives
\[
 -d''f_k'=
 \sum_{e=(i,j)}\left[
 d''\Phi_i\wedge d''\Phi_j\,h_{k-1}(C_i,C_j)
 +\frac12d''c_e\wedge
    (d''\Phi_iC_i^{k-1}+d''\Phi_jC_j^{k-1})\right].
 \tag{CF10}
\]
Use $h=-\Psi$ in the actual full MT15 formula.  Applying $d''$
to its first component, and remembering that CS36 subtracts
the full MT triple, gives the second through fourth lines of
(CF9).  The only polynomial simplifications are the finite
convolutions
\[
 \sum_{k=1}^{s-1}h_{s-k-1}(R_*,R)h_{k-1}(C_i,C_j)
                       =h_{s-2}(R_*,R,C_i,C_j)
 \tag{CF11}
\]
and the analogous identity with one $C_i$ argument.
Finally the scalar correction in CS36 has no first component.
For its independently closed-pair correction, the first component
is $d''h_\Psi\wedge(a_\nu d''b_\nu)$, where
$h_\Psi=-\Psi h_\ell(R_*,R)$.  Its second-colour derivative is
$d''\Psi\wedge d''a_\nu\wedge d''b_\nu h_\ell(R_*,R)$.
This gives the final line and proves the complete formula.
\end{proof}

Summing (CF9) with the six GR12 signs computes the raw cube
curl.  By GR16 it differs from (CF4) by the corresponding
$D$-boundary of the first quadruple component; it is not
legitimate to substitute one curl for the other without that
term.

\subsection{A full native facet where the proposed literal identity fails}

Fix $r=4$.  Choose $\varpi\in K$, $0<|\varpi|<1$, and put
$\lambda=v(\varpi)>0$.  On the standard angular torus in
$\mathbf P^3_K$, write
\[
 x_i=v(z_i/z_4)\quad(1\le i\le3),\qquad x_4=0.
 \tag{CF12}
\]
Take three independent original-$K$ torus units with valuations
$a,b,c$, multiplying entries one, two and three respectively.
Use the fixed positive cap constant $\varepsilon=4\lambda$.
Let $O$ be the ordinary valuation-box open of radius
$\lambda/10$ about
\[
 (x_1,x_2,x_3;a,b,c)
                  =(-2,-2,8;3,8,7)\lambda.
 \tag{CF13}
\]
The displayed center may lie on a facet; the open box contains
both sides.  All three weights are positive on $O$.

In the affine frame $z_4$, put $u=\min(x_1,x_2)$.
On the \emph{whole} open $O$, the eight full tuple minima are
\[
\begin{array}{c|cccccccc}
 v&0&A&B&AB&C&AC&BC&ABC\\\hline
 \phi_v&u&x_2&x_1&0&u&x_2&x_1&0.
\end{array}
 \tag{CF14}
\]
The inequalities follow directly from (CF13): the third entry
is strictly inactive, the first raised entry exceeds zero,
the second raised entry exceeds zero, and the unchanged fourth
entry has valuation zero.  These are strict inequalities on
the whole box, not identities inferred from rational points.

The raw face caps, in the exact order GR12, are likewise
\[
\begin{array}{c|cccccc}
 \nu&1&2&3&4&5&6\\\hline
 (a_\nu,b_\nu;\text{base})
       &(a,b;0)&(a,c;B)&(b,c;0)&(b,c;A)&(a,c;0)&(a,b;C)\\
 \sigma_\nu&+1&+1&+1&-1&-1&-1\\
 \Psi_\nu^{\rm raw}&0&0&x_1&0&x_2&0.
\end{array}
 \tag{CF15}
\]
Indeed each passive minimum is the displayed affine function,
and the alternative full-minimum-plus-weight-plus-cap branch
is strictly larger on $O$.  This verifies the entire caps,
not just their values on the center facet.

Now work in the standard projective metric: its reference
curvature is $C_0=-Du$ on this chart, its vertex scalars are
$\theta_v=\phi_v-u$, and its cap scalars are
$\kappa_\nu=\Psi_\nu^{\rm raw}-u$.
The only corner carrier meeting $O$ is
\[
 H=\{x_1=x_2\},\qquad C_0=\omega_H,
 \tag{CF16}
\]
with native weight one and positive standard normalization.
All $C_v$ vanish on $O$ except $C_0$ and $C_C=C_0$.
Every raw cap in (CF15) is affine, so
$C_{\Psi_\nu}=0$ for all six faces.  On $H$ write
$dx=d''x_1|_H=d''x_2|_H$.  The full tangential data are
\[
\begin{array}{c|cccccccc}
 v&0&A&B&AB&C&AC&BC&ABC\\\hline
 d''\theta_v|_H&0&0&0&-dx&0&0&0&-dx
\end{array}
 \tag{CF17}
\]
and
\[
 (d''\kappa_1,\ldots,d''\kappa_6)|_H
                         =(-dx,-dx,0,-dx,0,-dx).
 \tag{CF18}
\]

Every edge term of (CF3), after multiplication by
$d''\kappa_\nu$ and restriction to $H$, vanishes: two of
its one-form factors are zero or multiples of $dx$.
Using (CF8), only the final unit-pair term remains.  Its
six signed contributions give
\[
 \left.d''(E_{\rm glob})_A\right|_O
       =C_0\wedge dx\wedge(d''b-d''a)\wedge d''c.
 \tag{CF19}
\]
Here restriction to $H$ denotes the coefficient on the
$C_0$ carrier; the right side is the indicated actual
pure delta form on $O$.

The independently closed unit triple CU6 has first component
$-a\,d''b\wedge d''c$.  Therefore its expected comparison
curl in this rank is
\[
                d''(C_0E_{abc})_A
                   =-C_0\wedge d''a\wedge d''b\wedge d''c.
 \tag{CF20}
\]
Subtracting (CF20) from (CF19) gives the nonzero actual form
\[
 \boxed{\quad
 d''\bigl(E_{\rm glob}-C_0E_{abc}\bigr)_A
 =C_0\wedge
   \left[dx\wedge(d''b-d''a)\wedge d''c
                         +d''a\wedge d''b\wedge d''c\right]
                         \ne0\quad\text{on }O.
 \quad}
 \tag{CF21}
\]
For an explicit nonvanishing check, the bracket restricted to
the tangent vectors in the $a,b,c$ directions, with $x$ fixed,
has coefficient one.  These directions are independent on
$H$.  The identity torus chart has weight one, so this is a
nonzero native actual carrier coefficient.  Equivalently,
wedging with an ordinary smooth compact complementary form
supported near a relative interior point of $H$ gives a
nonzero integral.  The assertion is thus not merely a
many-to-one tropical pushforward calculation.

The same small certificate also rejects the analogous literal
curl identity for the raw cube in the standard projective
metric.  In (CF9), all normalized endpoint and cap derivatives
on $H$ are multiples of $dx$.  Its first line and all its
edge-dependent metric terms vanish there.  Its last line
has $h_1(C_{\Psi_\nu},C_0)=C_0$, so the same computation
gives (CF19) for the raw first curl.  This assertion is only
for this supplied metric and carrier; it does not identify
the two full curls in general or discard their $D$-difference.

\paragraph{The precise rejected step.}
Equation (CF21) disproves the proposed general identity
$d''(E_{\rm glob}-C_0^{r-3}E_{abc})_A=0$.
It therefore blocks applying a curl-free conical solver
directly to this uncorrected GR remainder.  It does
\emph{not} prove that the nonzero curl has no actual
$D$-primitive.  A prior explicit $D$-correction could remove
it, and a different genuinely extending error is also
permitted.  Neither possibility is decided here.

\paragraph{Actual interfaces and exact dependencies.}
The formulas (CF1)--(CF11) are finite actual differential
polynomials on the original $W$, with all full mixed
coefficients and every unit-weight wall retained.  Their
construction commutes with the actual sheaf restrictions
and all allowed refinements; arbitrary relations among
the supplied analytic units are preserved.  The example
uses an original-$K$ faithful native torus only to disprove
a particular universal identity.  It does not replace
the actual topology or test class of a general theorem.

The exact frozen-in-R11 GR source is
\path{c47cbd11e549bf94430182888bccffce20cad14ddc70f2929cab7ba04f5d3997}.
The full edge formula CS has hash
\path{d0c06d21c2289d292fb735d6966d456b9e05f1e4fab10f263a5007169d79cacf};
the complete MT source has hash
\path{d1365f528071a8db5e20f55628c88b84a06e73965dd808e8b27375c776fc73d4}.
The closed-pair CPT interface has hash
\path{725be5ed2460e1a0bbac88275b16533d5b5e8cc34a9befb3a6074c11f8e555f6},
and the unit triple CU6 has hash
\path{d38d0c77f879e3ca80962fa0fac2e56f6ea159fd59971a0bd810fbadc0a91471}.
No missing recovery module, full-kernel hypothesis, auxiliary
probe topology, arbitrary form descent or other branch is used.
